% Generated draft: requires mathematical and source review.
\documentclass[11pt]{article}
\usepackage[margin=25mm]{geometry}
\usepackage{fontspec}
\setmainfont{Latin Modern Roman}
\usepackage{amsmath,amssymb,mathtools,unicode-math}
\setmathfont{Latin Modern Math}
\usepackage{enumitem,needspace,xurl,hyperref}
\setlength{\parindent}{0pt}
\setlength{\parskip}{5pt}
\setlength{\emergencystretch}{3em}
\newcommand{\N}{\mathbb{N}}
\newcommand{\Z}{\mathbb{Z}}
\newcommand{\Q}{\mathbb{Q}}
\newcommand{\R}{\mathbb{R}}
\newcommand{\C}{\mathbb{C}}
\newcommand{\F}{\mathbb{F}}
\newcommand{\A}{\mathbb{A}}
\newcommand{\PP}{\mathbb{P}}
\newcommand{\E}{\mathbb{E}}
\newcommand{\eps}{\varepsilon}
\newcommand{\dd}{\,\mathrm{d}}
\begin{document}
\section*{UnsolvedMath ID: 30002523}
\textbf{Source problem code:} catalogue entry
\subsection{Definitions and mathematical statement}
Write \(k\) for the ground field and \(A\) for an associative unital finite-dimensional \(k\)-algebra. The supplied source does not specify the ground field or any restrictions on it; the formulation below is relative to a fixed ground field, whose intended hypotheses require confirmation. Use the convention that \(A\) is Gorenstein when its regular left and right modules both have finite injective dimension:
\[
\operatorname{injdim}_A({}_A A)<\infty,
\qquad
\operatorname{injdim}_{A^{\mathrm{op}}}(A_A)<\infty.
\]
Let \(S_1,\ldots,S_r\) represent the isomorphism classes of simple left \(A\)-modules, and put \(S=\bigoplus_{i=1}^r S_i\). Set
\[
A^e=A\otimes_k A^{\mathrm{op}},\qquad
\operatorname{HH}^{*}(A)=\bigoplus_{n\geq 0}\operatorname{Ext}^{n}_{A^e}(A,A),
\qquad
E_A(S)=\bigoplus_{n\geq 0}\operatorname{Ext}^{n}_A(S,S).
\]
Here \(\operatorname{HH}^{*}(A)\) is the Hochschild cohomology ring with its cup product, and \(E_A(S)\) is the Yoneda algebra with its Yoneda product. Equip \(E_A(S)\) with the usual graded \(\operatorname{HH}^{*}(A)\)-module structure induced by the Hochschild cohomology action on extensions. Condition \(\mathrm{Fg}\) means that \(\operatorname{HH}^{*}(A)\) is a noetherian ring and that \(E_A(S)\) is finitely generated as a graded module over this ring.

\textbf{Problem.} Characterize the finite-dimensional Gorenstein algebras \(A\) for which \(\mathrm{Fg}\) holds; equivalently, determine necessary and sufficient conditions, within this class, for
\[
\operatorname{HH}^{*}(A)\text{ to be noetherian}
\quad\text{and}\quad
E_A(S)\text{ to be finitely generated over }\operatorname{HH}^{*}(A).
\]
The source asks for a characterization, not an assertion that every Gorenstein algebra satisfies \(\mathrm{Fg}\). The choice of left modules and the two-sided injective-dimension convention make the terminology explicit; the source does not itself specify these conventions or the intended ground-field hypotheses.
\subsection{Short English statement}
Determine which finite-dimensional Gorenstein algebras have noetherian Hochschild cohomology and a Yoneda algebra of the direct sum of representatives of the simple modules that is finitely generated over Hochschild cohomology. The intended assumptions on the ground field remain unspecified. The supplied registry labels the problem partially solved; that label and its dated literature assessment are not independently verified here.
\subsection{Sources}
\begin{itemize}
\item Local UnsolvedMath registry, permanent ID \(30002523\), title \emph{Finite Generation for Gorenstein Algebras}, source problem code \texttt{catalogue entry}, in the Oberwolfach Reports Open Problems collection. The supplied record attributes the problem to the 2014 report but names no individual proposer.
\item Supplied primary reference: \emph{Representation Theory of Quivers and Finite Dimensional Algebras} (2014), Oberwolfach Reports, pp.\ 453--529, DOI \texttt{10.4171/owr/2014/08}. Supplied literature-source label: Oberwolfach Report 08/2014. \url{https://doi.org/10.4171/owr/2014/08}.
\item Supplied literature source: \emph{Finite generation for Gorenstein monomial algebras}. \url{https://arxiv.org/abs/1909.00487}.
\item Supplied literature source: \emph{Combinatorial characterisation for monomial algebras}. \url{https://arxiv.org/abs/2301.05476}.
\item Literature assessment dated \texttt{2026-08-21}: the supplied curation states that finite generation implies Gorensteinness, that the converse is proved for finite-dimensional monomial algebras, and that further special classes have been characterized, while the general characterization remains open. This is supplied dated curation, not an independently verified assessment; the record expressly requests independent verification of time-sensitive status and cited claims.
\end{itemize}
\end{document}
